% ------------------------------------------------------------------------------
% 4.1. INTRODUCCIÓN
% ------------------------------------------------------------------------------
\subsection{4.1. Introducción}\label{sec:ieee_introduccion}

\subsubsection{4.1.1. Propósito}\label{sec:ieee_proposito}
% [INSTRUCCIÓN]: Definir el propósito específico del documento ERS y su audiencia objetivo.

El presente documento tiene como propósito definir, de manera clara, completa y sin ambigüedades, las funcionalidades y restricciones requeridas para el desarrollo del Sistema de Información Web y Móvil para la Postulación y Preselección Automatizada de la Beca Comedor mediante Inteligencia Artificial, con Reportes de Business Intelligence para la DUBSS - UAGRM. Esta Especificación de Requisitos del Software (ERS) está dirigida al equipo de desarrollo encargado de la construcción del sistema, al personal directivo y administrativo de la Dirección Universitaria de Bienestar Social y Salud (DUBSS) y a los docentes evaluadores del proyecto de grado, con el objetivo de establecer un acuerdo formal sobre el alcance, el comportamiento esperado y las restricciones del software antes de iniciar las etapas de diseño e implementación.

\subsubsection{4.1.2. Ámbito del Sistema}\label{sec:ieee_ambito}
% [INSTRUCCIÓN]: Nombre del producto software, alcance operativo, beneficios y metas generales.

\textbf{Nombre del proyecto:} Sistema de Información Web y Móvil para la Postulación y Preselección Automatizada de la Beca Comedor mediante Inteligencia Artificial, con Reportes de Business Intelligence para la DUBSS - UAGRM.

El sistema comprenderá las siguientes funcionalidades principales:

\begin{itemize}
    \item Registro y autenticación de postulantes a través de las plataformas web y móvil.
    \item Postulación en línea a la Beca Comedor mediante el llenado de formularios socioeconómicos digitales.
    \item Preselección automatizada de solicitudes empleando un modelo de Inteligencia Artificial y Machine Learning que precalifica algorítmicamente el nivel de vulnerabilidad socioeconómica de cada postulante.
    \item Gestión administrativa de los expedientes finalistas, incluyendo la auditoría física de respaldos y el registro de los estudiantes efectivamente becados.
    \item Registro y control de la asistencia física y digital de los becarios al Comedor Universitario.
    \item Generación de reportes y dashboards de Business Intelligence con proyecciones predictivas que apoyen la toma de decisiones institucionales.
\end{itemize}

Entre los principales beneficios que otorgará el sistema a la DUBSS se encuentran la reducción drástica del tiempo y del personal destinado a la revisión manual de carpetas, la eliminación del cuello de botella administrativo generado por las 20,000 solicitudes recibidas por gestión, la incorporación de un criterio de selección objetivo, transparente y auditable, y la disposición de información oportuna y confiable para la planificación y rendición de cuentas institucional.

\subsubsection{4.1.3. Definiciones, Acrónimos y Abreviaturas}\label{sec:ieee_definiciones}
% [INSTRUCCIÓN]: Glosario de términos técnicos, siglas y acrónimos necesarios para interpretar el documento.

La Tabla~\ref{tab:definiciones_abreviaturas} resume los principales términos, siglas y acrónimos empleados a lo largo del presente documento.

\begin{table}[H]
\caption{Definiciones, acrónimos y abreviaturas empleados en el documento}
\label{tab:definiciones_abreviaturas}
\begin{threeparttable}
\begin{tabular}{>{\raggedright\arraybackslash}p{3cm} >{\raggedright\arraybackslash}p{10cm}}
\toprule
Nombre & Descripción \\
\midrule
DUBSS & Dirección Universitaria de Bienestar Social y Salud, unidad de la UAGRM responsable de la Beca Comedor. \\
Administrador & Personal de la DUBSS que audita, valida y gestiona las postulaciones dentro del sistema. \\
Postulante & Estudiante de la UAGRM que solicita el beneficio de la Beca Comedor a través del sistema. \\
Usuario & Persona que interactúa con el software, ya sea en calidad de postulante o de administrador. \\
RF & Requerimiento funcional. \\
RNF & Requerimiento no funcional. \\
ERS & Especificación de Requisitos del Software. \\
PCN & Prueba de caja negra. \\
IA & Inteligencia Artificial. \\
ML & Machine Learning (aprendizaje automático). \\
PLN & Procesamiento de Lenguaje Natural. \\
API & Interfaz de Programación de Aplicaciones. \\
UI/UX & Interfaz de Usuario y Experiencia de Usuario. \\
JS & Lenguaje de programación JavaScript. \\
\bottomrule
\end{tabular}
\begin{tablenotes}
\item \textit{Nota.} Elaboración propia.
\end{tablenotes}
\end{threeparttable}
\end{table}

\subsubsection{4.1.4. Referencias}\label{sec:ieee_referencias}
% [INSTRUCCIÓN]: Lista de documentos de referencia, normativas institucionales y estándares técnicos.

Los principales estándares, libros y artículos empleados para las definiciones y el desarrollo de este capítulo son los siguientes:

\begin{itemize}
    \item \textcite{ieee830_1998} y \textcite{iso29148_2011}.
    \item \textcite{pressman2010ingenieria}.
    \item \textcite{sommerville2011ingenieria}.
\end{itemize}

\subsubsection{4.1.5. Visión General del Documento}\label{sec:ieee_vision_general}
% [INSTRUCCIÓN]: Breve descripción de la organización del resto de la especificación de requisitos.

El presente documento se compone de tres secciones (Introducción, Descripción General y Requisitos Específicos) que describen detalladamente el producto a desarrollarse. En primer lugar, la sección de Introducción establece el propósito, el ámbito, las definiciones y la visión general del documento. A continuación, la sección de Descripción General sitúa al sistema dentro de su contexto de uso, describiendo sus funciones principales, las características de sus usuarios y las restricciones y suposiciones bajo las cuales operará. Finalmente, la sección de Requisitos Específicos detalla de manera formal cada uno de los requisitos funcionales, no funcionales y de interfaz de usuario que deberá satisfacer el sistema.
